\chapter*{Como ler este livro}
\addcontentsline{toc}{chapter}{Como ler este livro}
\markboth{Como ler este livro}{Como ler este livro}

\section*{A arquitetura do argumento}

O livro está organizado em cinco partes que correspondem a cinco momentos de uma
mesma investigação, e a estrutura cabe em um parágrafo.
A Parte~I ensina o necessário. A Parte~II apresenta a base --- três
estratégias que, juntas, constituem o primeiro sistema completo de aprendizado de
mundo aberto para poços de petróleo. A Parte~III mostra como essa base foi revista
quando suas limitações ficaram claras. A Parte~IV acrescenta a camada que faltava:
explicação. A Parte~V amarra tudo e discute o que falta.

\section*{Trilhas de leitura}

\begin{description}[leftmargin=0pt,style=nextline,itemsep=0.7em]

\item[Leitura completa (curso de um semestre).]
Capítulos 1 a 18 na ordem. Os \cref{cap:introducao,cap:dominio,cap:dados} cobrem
domínio e dados; \cref{cap:aprendizado,cap:profundo,cap:novidades} cobrem método;
os demais são aplicação. Uma disciplina de quinze semanas comporta um capítulo por
semana com três semanas de folga para os projetos.

\item[Leitura para quem já sabe aprendizado de máquina.]
Comece pelo \cref{cap:dominio} (o vocabulário de poços é indispensável e não é óbvio),
leia o \cref{cap:dados} (o conjunto 3W tem particularidades que confundem quem chega
de outros domínios), passe rapidamente pelo \cref{cap:novidades} e vá para a
\cref{parte:base}.

\item[Leitura para quem quer implantar.]
\cref{cap:dominio}, \cref{cap:dual} (o sistema que efetivamente está em campo),
\cref{cap:sintese} (arquitetura de referência) e \cref{cap:operacao} (implantação,
deriva, governança). Volte aos demais conforme a necessidade.

\item[Leitura para quem só quer os resultados.]
\cref{cap:dual}, \cref{cap:binarios}, \cref{cap:mundoaberto},
\cref{cap:segmentacao}, \cref{cap:mahalanobis} e \cref{cap:resultados-agente}
contêm todas as tabelas. O \cref{cap:sintese} consolida.

\end{description}

\section*{Convenções tipográficas}

Ao longo do texto você encontrará seis elementos recorrentes.

\begin{objetivos}
Abre cada capítulo. Lista, em forma de verbos de ação, o que se espera que o leitor
saiba fazer ao terminá-lo. Os verbos foram escolhidos para refletir o que o texto de
fato entrega: este é um livro de argumentação e evidência, não um manual de
implementação. Onde o objetivo é ``descrever'' ou ``reconhecer'', o capítulo não
pretende habilitar o leitor a implementar o método do zero; para isso, o
\cref{cap:pratica} aponta código executável e o \cref{cap:reproducao} reúne os
hiperparâmetros.
\end{objetivos}

\medskip
\noindent
As \textbf{caixas de definição}, em azul, introduzem conceitos formais e são numeradas
por capítulo (por exemplo, \emph{Definição 6.2}), de modo que podem ser referenciadas
cruzadamente. As \textbf{caixas de exemplo}, em verde, aplicam o conceito imediatamente
anterior a um caso concreto, quase sempre extraído dos dados reais discutidos no livro,
e seguem a mesma numeração.

\begin{resultadochave}
Destaca um número ou achado que sustenta o argumento do capítulo. É o que você deve
levar consigo se ler apenas as caixas.
\end{resultadochave}

\begin{destaque}[Caixa de atenção]
Marca uma armadilha conceitual, uma limitação relevante ou uma decisão de projeto cuja
motivação não é evidente.
\end{destaque}

\medskip
\noindent
Cada capítulo termina com um \textbf{resumo}, uma lista de \textbf{exercícios} ---
mesclando questões conceituais, exercícios de cálculo e propostas de implementação ---
e sugestões de \textbf{leituras recomendadas}.


\section*{Nomenclatura das classes de anomalia}

Todo o livro utiliza a numeração canônica do conjunto de dados 3W, apresentada em
detalhe no \cref{cap:dados}. Por conveniência, ela é reproduzida aqui:

\begin{center}
\small
\begin{tabular}{cl}
\toprule
\textbf{Classe} & \textbf{Evento} \\
\midrule
0 & Estado normal \\
1 & Aumento abrupto de BSW \\
2 & Fechamento espúrio da DHSV \\
3 & Golfadas severas \\
4 & Instabilidade de fluxo \\
5 & Perda rápida de produtividade \\
6 & Restrição rápida no PCK \\
7 & Incrustação no PCK \\
8 & Hidrato na linha de produção \\
9 & Hidrato na linha de serviço \\
\midrule
--- & Falha de ICV \emph{(dado proprietário; não pertence ao 3W)} \\
--- & Operação de válvula \emph{(dado proprietário; não pertence ao 3W)} \\
\bottomrule
\end{tabular}
\end{center}

\section*{Sobre reprodutibilidade}

O conjunto de dados 3W é público. Os dados proprietários de poços da Petrobras
utilizados nos estudos de caso não o são, e portanto os resultados que dependem
exclusivamente deles não podem ser reproduzidos por terceiros. Sempre que um resultado
depende de dados fechados, isso é dito explicitamente. O \cref{cap:reproducao} reúne
hiperparâmetros, ambiente computacional e procedimentos de avaliação de todos os
experimentos discutidos.

O leitor que quiser executar código, e não apenas lê-lo, deve começar pelo
\cref{cap:pratica}. Ele descreve um recurso didático público, hospedado no próprio
repositório do 3W \citep{Petrobras3W,Lopes2026Intro3W}, que percorre em cinco cadernos
Jupyter todo o caminho da \cref{parte:fundamentos}: tratamento de dados, visualização,
autoencoders, classificadores supervisionados e agrupamento. As caixas de
\textbf{prática} espalhadas pela \cref{parte:fundamentos} apontam, capítulo a capítulo,
qual caderno exercita qual conceito.
